\lysh{38930}{4}

\begin{alist}
\item
\begin{rlist}
\item O όγκος του σφαιρικού τμήματος ύψους $h$, που περιέχει το απορρυπαντικό, είναι ίσος
με το $\dfrac{1}{3}$ του όγκου της κάψουλας, οπότε:
$\dfrac{\pi h^2}{3} (3r - h) = \dfrac{1}{3} \cdot \left( \dfrac{4}{3} \pi r^3 \right)$, οπότε για $r = 2\text{ cm}$ έχουμε
\begin{gather*}
\dfrac{h^2}{3} (6 - h) = \dfrac{1}{3} \cdot \left( \dfrac{4}{3} \cdot 8 \right)\text{, δηλαδή} \\
\dfrac{6h^2 - h^3}{3} = \dfrac{32}{9}\text{και τελικά}
\end{gather*}
$3h^3 - 18h^2 + 32 = 0$. (1)

\item Η εξίσωση που βρήκαμε στο $\alpha)$ii. ερώτημα είναι ισοδύναμη με την $h^3 - 6h^2 + \dfrac{32}{3} = 0$.
Οι λύσεις της εξίσωσης αυτής είναι οι τετμημένες των κοινών σημείων της γραφικής παράστασης της συνάρτησης $f(x) = x^3 - 6x^2 + \dfrac{32}{3}$ με τον άξονα $x'x$. Από το σχήμα βλέπουμε ότι η γραφική παράσταση της συνάρτησης έχει τρία κοινά σημεία με τον
άξονα $x'x$, όμως η τετμημένη μόνο του ενός θα μπορούσε να είναι λύση της εξίσωσης (1) (το $h$, ως ύψος σφαιρικού τμήματος δεν μπορεί να είναι αρνητικός αριθμός και δεν μπορεί να είναι μεγαλύτερο από την ακτίνα $r = 2\text{ cm}$ της σφαίρας). Άρα το ύψος $h$ του
σφαιρικού τμήματος που περιέχει το απορρυπαντικό είναι περίπου $h = 1,5\text{ cm}$.
\end{rlist}

\item Τα $\dfrac{2}{9}$ ενός κυλίνδρου ίδιας ακτίνας ($\rho = 2\text{ cm}$) και ύψους $h = 4\text{ cm}$ (όσο η διάμετρος
της σφαιρικής κάψουλας) έχουν όγκο:
\[V_{1} = \dfrac{2}{9} \cdot (\pi \cdot 2^2 \cdot 4) = \dfrac{32\pi}{9}\text{ cm}^3.\]

Ο όγκος του απορρυπαντικού είναι το $\dfrac{1}{3}$ του όγκου της σφαίρας ακτίνας $r = 2\text{ cm}$, άρα
\[V_{2} = \dfrac{1}{3} \cdot \dfrac{4}{3} \pi \cdot 2^3 = \dfrac{32\pi}{9}\text{ cm}^3.\]

Άρα ο ισχυρισμός είναι σωστός.
\end{alist}
